\begin{pastexam}[label={px:ntu-107-dsa-05}]{臺大資工所 107 學年度 第 V 題}
（問答）

(35 points) In the data structure \textit{queue}, elements can be added to one end and removed
from the other. Different from queues, a \textit{deque} allows elements to be added to and removed
from both ends.
\begin{enumerate}[label=(\alph*)]
\item (15 points) Given a deque with a permutation of the numbers from 1 to $n$ ($n\ge5$), we want
to remove all of the numbers to form a sequence $<a_1,a_2,a_3,\ldots,a_n>$ that maximizes the
alternative sum $(a_1-a_2+a_3\ldots\pm a_n)$, where $a_1$ is the first number removed, $a_2$ is the
second number, and so on. For example, suppose a deque contains 5 numbers as follows:

\hspace*{2em}(head) 2, 3, 1, 4, 5 (rear)

One possible generated sequence is $<5,2,4,1,3>$ whose alternative sum is 9.
\begin{enumerate}[label=(\arabic*)]
\item (5 points) Give an example to demonstrate that the problem cannot be solved by a greedy
algorithm, which first removes the larger one of the two ends, then removing in the order of the
smaller one, the larger one, the smaller one, etc. Show your calculation.
\item (10 points) Devise an $O(n^2)$ algorithm to generate the optimal sequence for a given deque.
Explain how your algorithm works on the example you give previously.
\end{enumerate}
\item (20 points) Stacks can be used to implement queues and deques.
\begin{enumerate}[label=(\arabic*)]
\item (6 points) Describe how to implement a queue using only two stacks. And then describe how to
extend the way to implement a queue to implement a deque using the same two stacks.
\item (7 points) Use the accounting method to prove the amortized time complexity of a sequence of
$n$ queue operations (implemented with stacks) is $O(n)$. Show your derivation.
\item (7 points) Prove or disprove that the amortized time complexity of a sequence of $n$ deque
operations (implemented with stacks) is also $O(n)$. Show your derivation.
\end{enumerate}
\end{enumerate}
\end{pastexam}

\begin{pastsoln}
本題橫跨幾章：(a) 是動態規劃（資結書第~18~章），(b) 是用堆疊做佇列（資結書第~6~章）和記帳法
（資結書第~3.2~節）。 (b) 佔的分數較多，歸在平攤分析這一章。

\textbf{(a)(1) 貪心的反例。} 取 $n = 5$，(head) $3, 1, 2, 4, 5$ (rear)。 照題目的貪心法：第奇數次拿兩端
較大的（要加），第偶數次拿兩端較小的（要減）。
\begin{center}
\small
\begin{tabular}{clcl}
\toprule
第幾次 & 兩端 & 拿走 & 剩下 \\
\midrule
1（加，拿大的） & $3$、$5$ & $5$ & $3, 1, 2, 4$ \\
2（減，拿小的） & $3$、$4$ & $3$ & $1, 2, 4$ \\
3（加，拿大的） & $1$、$4$ & $4$ & $1, 2$ \\
4（減，拿小的） & $1$、$2$ & $1$ & $2$ \\
5（加） & $2$ & $2$ & （空） \\
\bottomrule
\end{tabular}
\end{center}
得到 $\langle 5, 3, 4, 1, 2\rangle$，交錯和 $5 - 3 + 4 - 1 + 2 = 7$。 但是先從頭拿 $3$、再從頭拿 $1$、
然後從尾拿 $5$、從頭拿 $2$、最後拿 $4$，得到 $\langle 3, 1, 5, 2, 4\rangle$，交錯和
$3 - 1 + 5 - 2 + 4 = 9$。 $9$ 已經是 $n = 5$ 的上限：交錯和是三個「加」的數減兩個「減」的數，
最多 $(5 + 4 + 3) - (2 + 1) = 9$。 所以貪心法得不到最佳解。 它錯在第 1 步：拿走 $5$ 雖然當下最
大，卻讓第 2 步只能在 $3$、$4$ 裡挑一個減掉，大的數被浪費在「減」的位置；先忍住、把 $3$ 加進
來、把 $1$ 減掉，$5$ 和 $4$ 之後都還拿得到。

\textbf{(a)(2) 動態規劃。} 不管之前怎麼拿，deque 裡剩下的永遠是原序列連續的一段 $d_i, \ldots, d_j$
（$d_1, \ldots, d_n$ 是從 head 到 rear 的原序列）。 已經拿走 $n - (j - i + 1)$ 個，所以下一個拿走的數
在序列中的位置、也就是它要加還是減，完全由這一段的長度決定：令
\[
s(L) = \begin{cases} +1 & n - L \text{ 是偶數（下一個是第奇數個）},\\ -1 & n - L \text{ 是奇數}.\end{cases}
\]
以後怎麼拿只和這一段有關，和之前的順序無關，所以可以定義子問題（資結書第~18.1.5~節的步驟）：
$D[i, j]$ 是「剩下 $d_i..d_j$ 時，把它們拿完能得到的最大加減總和」。 下一個不是拿頭 $d_i$ 就是拿尾
$d_j$：
\[
D[i, i] = s(1)\, d_i, \qquad
D[i, j] = \max\bigl(s(L)\, d_i + D[i+1, j],\; s(L)\, d_j + D[i, j-1]\bigr),\quad L = j - i + 1.
\]
答案是 $D[1, n]$。 依長度 $L = 1, 2, \ldots, n$ 填表，同時記下每一格選了頭還是尾，最後從 $[1, n]$ 照
記錄走回來，就得到最佳的拿取順序。
\begin{pseudo}{雙端佇列交錯和的最大值}
\Function{Max-Alternating}{$d, n$}
    \For{$L \gets 1$ \To $n$}
        \State $s \gets +1$ \textbf{if} $n - L$ 是偶數 \textbf{else} $-1$
        \For{$i \gets 1$ \To $n - L + 1$}
            \State $j \gets i + L - 1$
            \If{$L = 1$}
                \State $D[i, i] \gets s \cdot d_i$；$\mathit{take}[i, i] \gets \text{頭}$
            \Else
                \State $x \gets s \cdot d_i + D[i+1, j]$；$y \gets s \cdot d_j + D[i, j-1]$
                \State $D[i, j] \gets \max(x, y)$；$\mathit{take}[i, j] \gets$（$x \ge y$ 時頭，否則尾）
            \EndIf
        \EndFor
    \EndFor
    \State $i \gets 1$；$j \gets n$
    \While{$i \le j$} \Comment{照記錄輸出順序}
        \If{$\mathit{take}[i, j] = \text{頭}$}
            \State 輸出 $d_i$；$i \gets i + 1$
        \Else
            \State 輸出 $d_j$；$j \gets j - 1$
        \EndIf
    \EndWhile
    \State \Return $D[1, n]$
\EndFunction
\end{pseudo}
正確性：最佳的拿法第一步拿頭或拿尾；拿頭之後，剩下 $d_{i+1}..d_j$ 的拿法如果不是那一段的最佳，換
成最佳的只會讓總和變大，所以 $D[i, j]$ 一定等於兩個選項之一，遞迴式成立（最佳子結構）。 時間：表有
$n(n+1)/2$ 格，每格 $O(1)$，回溯 $O(n)$，總共 $O(n^2)$；空間 $O(n^2)$。

\textbf{在 (a)(1) 的例子上跑一次}：$d = \langle 3, 1, 2, 4, 5\rangle$，$n = 5$。 長度 $1, 3, 5$ 的段
$s = +1$，長度 $2, 4$ 的段 $s = -1$。
\begin{center}
\small
\begin{tabular}{cll}
\toprule
$L$ & 計算（頭；尾） & 結果 \\
\midrule
1 & $D[1,1] = 3$，$D[2,2] = 1$，$D[3,3] = 2$，$D[4,4] = 4$，$D[5,5] = 5$ & \\
2 & $D[1,2] = \max(-3 + 1,\ -1 + 3) = 2$ & 尾 \\
 & $D[2,3] = \max(-1 + 2,\ -2 + 1) = 1$ & 頭 \\
 & $D[3,4] = \max(-2 + 4,\ -4 + 2) = 2$ & 頭 \\
 & $D[4,5] = \max(-4 + 5,\ -5 + 4) = 1$ & 頭 \\
3 & $D[1,3] = \max(3 + D[2,3],\ 2 + D[1,2]) = \max(4, 4) = 4$ & 頭 \\
 & $D[2,4] = \max(1 + D[3,4],\ 4 + D[2,3]) = \max(3, 5) = 5$ & 尾 \\
 & $D[3,5] = \max(2 + D[4,5],\ 5 + D[3,4]) = \max(3, 7) = 7$ & 尾 \\
4 & $D[1,4] = \max(-3 + D[2,4],\ -4 + D[1,3]) = \max(2, 0) = 2$ & 頭 \\
 & $D[2,5] = \max(-1 + D[3,5],\ -5 + D[2,4]) = \max(6, 0) = 6$ & 頭 \\
5 & $D[1,5] = \max(3 + D[2,5],\ 5 + D[1,4]) = \max(9, 7) = 9$ & 頭 \\
\bottomrule
\end{tabular}
\end{center}
回溯：$[1,5]$ 拿頭 $3$ → $[2,5]$ 拿頭 $1$ → $[3,5]$ 拿尾 $5$ → $[3,4]$ 拿頭 $2$ → $[4,4]$ 拿 $4$，順序
$\langle 3, 1, 5, 2, 4\rangle$，交錯和 $9$，和 (a)(1) 一樣。 $D[1,5]$ 那一格的「尾」選項 $5 + D[1,4] = 7$
正是貪心法的第一步，表格顯示它比拿頭少 $2$。 （題目自己的例子 $2, 3, 1, 4, 5$ 用同樣的表算，最大值
也是 $9$，$\langle 5, 2, 4, 1, 3\rangle$ 就是一組最佳解。）

\textbf{(b)(1) 兩個堆疊做佇列、雙端佇列。} 兩個堆疊 $S_{\mathit{in}}$、$S_{\mathit{out}}$。
\begin{itemize}
\item \textbf{佇列}。 ENQUEUE$(x)$：push 到 $S_{\mathit{in}}$。 DEQUEUE：如果 $S_{\mathit{out}}$ 是空的，
把 $S_{\mathit{in}}$ 的元素一個一個 pop 出來 push 進 $S_{\mathit{out}}$（順序因此反過來，最早進來的到了
$S_{\mathit{out}}$ 頂端）；然後 pop $S_{\mathit{out}}$。 整個佇列由前到後是「$S_{\mathit{out}}$ 由頂到底，
接著 $S_{\mathit{in}}$ 由底到頂」，前端在 $S_{\mathit{out}}$ 的頂端，後端在 $S_{\mathit{in}}$ 的頂端。
\item \textbf{雙端佇列}。 沿用同樣的排法，$S_{\mathit{out}}$ 管前端、$S_{\mathit{in}}$ 管後端，兩邊對稱：
PUSH-FRONT$(x)$ push 到 $S_{\mathit{out}}$；PUSH-BACK$(x)$ push 到 $S_{\mathit{in}}$；POP-FRONT 同佇列的
DEQUEUE；POP-BACK：如果 $S_{\mathit{in}}$ 是空的，先把 $S_{\mathit{out}}$ 全部倒進 $S_{\mathit{in}}$，再 pop
$S_{\mathit{in}}$。 倒的時候順序反過來，正好讓原本在 $S_{\mathit{out}}$ 底部的最後一個元素到了
$S_{\mathit{in}}$ 頂端。
\end{itemize}

\textbf{(b)(2) 記帳法：佇列。} 成本模型：每次 push 或 pop 花 $1$。 收費訂成 ENQUEUE 收 $3$，DEQUEUE
收 $1$（資結書第~3.2~節）。
\begin{itemize}
\item ENQUEUE 實際花 $1$（push 到 $S_{\mathit{in}}$），多收的 $2$ 存在這個元素身上。
\item DEQUEUE 搬家時，每個被搬的元素花 $2$（從 $S_{\mathit{in}}$ pop、push 到 $S_{\mathit{out}}$），剛好
用它身上的 $2$ 付；最後 pop $S_{\mathit{out}}$ 的 $1$ 由 DEQUEUE 自己收的 $1$ 付。
\end{itemize}
存款不會變負：任何時刻，$S_{\mathit{in}}$ 裡每個元素身上剛好有 $2$，$S_{\mathit{out}}$ 裡的元素身上是
$0$，存款 $= 2\,|S_{\mathit{in}}| \ge 0$。 每個元素一生只會從 $S_{\mathit{in}}$ 搬到 $S_{\mathit{out}}$ 一次
（$S_{\mathit{out}}$ 的元素不會回到 $S_{\mathit{in}}$），所以它身上的 $2$ 夠用。 因此 $n$ 次操作的實際總
成本不超過總收費 $\le 3n$，是 $O(n)$，每次操作平攤 $O(1)$。

\textbf{(b)(3) 雙端佇列：不成立（反證）。} 用 (b)(1) 的做法，下面這串 $n$ 個操作要花 $\Theta(n^2)$。
令 $m = n/2$。 先 PUSH-BACK $m$ 次，元素全在 $S_{\mathit{in}}$。 接著 POP-FRONT、POP-BACK 交替做 $m$ 次：
\begin{itemize}
\item 第 $1$ 次 POP-FRONT：$S_{\mathit{out}}$ 空，把 $m$ 個全搬過去（$2m$），pop 一個；剩 $m - 1$ 個，全在
$S_{\mathit{out}}$。
\item 第 $1$ 次 POP-BACK：$S_{\mathit{in}}$ 空，把 $m - 1$ 個全搬回來（$2(m - 1)$），pop 一個；剩 $m - 2$ 個，
全在 $S_{\mathit{in}}$。
\item 以此類推，每一次都要把剩下的全部搬到另一邊。
\end{itemize}
總成本是
\[
m + 2\bigl(m + (m - 1) + \cdots + 1\bigr) + m = m(m + 1) + 2m = m^2 + 3m = \Theta(n^2).
\]
所以 $n$ 次操作的總時間不是 $O(n)$（平攤每次是 $\Theta(n)$，不是 $O(1)$），題目的敘述不成立。 (b)(2)
的論證在這裡失效，是因為元素可以在兩個堆疊之間來回搬很多次，身上的固定存款付不起。 （補充：如果
另外用一個暫存堆疊，一邊空了的時候只把另一邊的\textbf{一半}搬過來，可以證明平攤 $O(1)$；但那已經不是
「只用同樣兩個堆疊」的做法。）
\end{pastsoln}
